\begin{landscape}
\begingroup
\small
\renewcommand{\arraystretch}{1.15}
\setlength{\tabcolsep}{4pt}
\begin{longtable}{p{2.2cm}p{3.0cm}p{4.0cm}rccc}
\caption{Geographic location of missing observations in the individual-characteristics specifications}\label{tab:a10_missing_places}\\
\toprule
Region & Province & Municipality & N & Full & No blank/null vote & No ideology \\
\midrule
\endfirsthead
\multicolumn{7}{c}{\tablename\ \thetable{} -- continued from previous page} \\
\toprule
Region & Province & Municipality & N & Full & No blank/null vote & No ideology \\
\midrule
\endhead
\midrule
\multicolumn{7}{r}{Continued on next page} \\
\endfoot
\bottomrule
\endlastfoot
Pampeana/Centro & CHACO & CHACABUCO & 90 & 44 (48.9\%) & 31 (34.4\%) & 34 (37.8\%) \\
 &  & COMANDANTE FERNANDEZ & 108 & 43 (39.8\%) & 28 (25.9\%) & 27 (25.0\%) \\
 &  & GENERAL GUEMES & 73 & 28 (38.4\%) & 18 (24.7\%) & 21 (28.8\%) \\
 &  & LIBERTADOR GENERAL SAN MARTIN & 91 & 34 (37.4\%) & 21 (23.1\%) & 20 (22.0\%) \\
 &  & LIBERTAD & 90 & 33 (36.7\%) & 16 (17.8\%) & 22 (24.4\%) \\
 &  & SAN FERNANDO & 184 & 50 (27.2\%) & 23 (12.5\%) & 40 (21.7\%) \\
\addlinespace[2pt]
 & RIO NEGRO & GENERAL ROCA & 273 & 101 (37.0\%) & 59 (21.6\%) & 74 (27.1\%) \\
\addlinespace[2pt]
 & SALTA & ROSARIO DE LA FRONTERA & 91 & 31 (34.1\%) & 17 (18.7\%) & 20 (22.0\%) \\
 &  & CAPITAL & 90 & 30 (33.3\%) & 14 (15.6\%) & 25 (27.8\%) \\
 &  & LA CALDERA & 92 & 17 (18.5\%) & 5 (5.4\%) & 15 (16.3\%) \\
\addlinespace[2pt]
Region not identified & BUENOS AIRES & MAGDALENA & 90 & 43 (47.8\%) & 27 (30.0\%) & 26 (28.9\%) \\
 &  & LOMAS DE ZAMORA & 91 & 41 (45.1\%) & 24 (26.4\%) & 33 (36.3\%) \\
 &  & MORENO & 91 & 37 (40.7\%) & 18 (19.8\%) & 30 (33.0\%) \\
 &  & CANUELAS & 90 & 36 (40.0\%) & 22 (24.4\%) & 24 (26.7\%) \\
 &  & LUJAN & 180 & 72 (40.0\%) & 47 (26.1\%) & 53 (29.4\%) \\
 &  & VICENTE LOPEZ & 90 & 36 (40.0\%) & 23 (25.6\%) & 24 (26.7\%) \\
 &  & LA MATANZA & 453 & 175 (38.6\%) & 92 (20.3\%) & 127 (28.0\%) \\
 &  & MORON & 90 & 34 (37.8\%) & 16 (17.8\%) & 26 (28.9\%) \\
 &  & PERGAMINO & 90 & 34 (37.8\%) & 18 (20.0\%) & 21 (23.3\%) \\
 &  & LA PLATA & 91 & 32 (35.2\%) & 24 (26.4\%) & 21 (23.1\%) \\
 &  & SAN MIGUEL & 92 & 32 (34.8\%) & 18 (19.6\%) & 25 (27.2\%) \\
 &  & LANUS & 92 & 31 (33.7\%) & 14 (15.2\%) & 25 (27.2\%) \\
 &  & MERLO & 90 & 30 (33.3\%) & 18 (20.0\%) & 26 (28.9\%) \\
 &  & QUILMES & 18 & 6 (33.3\%) & 2 (11.1\%) & 5 (27.8\%) \\
 &  & BARADERO & 91 & 30 (33.0\%) & 22 (24.2\%) & 18 (19.8\%) \\
 &  & TIGRE & 252 & 83 (32.9\%) & 44 (17.5\%) & 62 (24.6\%) \\
 &  & BERISSO & 90 & 28 (31.1\%) & 17 (18.9\%) & 15 (16.7\%) \\
 &  & JOSE C. PAZ & 90 & 28 (31.1\%) & 17 (18.9\%) & 25 (27.8\%) \\
 &  & MAR CHIQUITA & 90 & 28 (31.1\%) & 19 (21.1\%) & 15 (16.7\%) \\
 &  & GENERAL PUEYRREDON & 183 & 55 (30.1\%) & 30 (16.4\%) & 39 (21.3\%) \\
 &  & ENSENADA & 91 & 27 (29.7\%) & 18 (19.8\%) & 17 (18.7\%) \\
 &  & BAHIA BLANCA & 90 & 26 (28.9\%) & 14 (15.6\%) & 20 (22.2\%) \\
 &  & AVELLANEDA & 185 & 52 (28.1\%) & 22 (11.9\%) & 47 (25.4\%) \\
 &  & SAN VICENTE & 91 & 24 (26.4\%) & 9 (9.9\%) & 18 (19.8\%) \\
\addlinespace[2pt]
 & CIUDAD AUTONOMA DE BUENOS AIRES & CITY OF BUENOS AIRES & 631 & 234 (37.1\%) & 115 (18.2\%) & 203 (32.2\%) \\
\addlinespace[2pt]
 & CORDOBA & PUNILLA & 90 & 49 (54.4\%) & 28 (31.1\%) & 43 (47.8\%) \\
 &  & SAN JUSTO & 90 & 41 (45.6\%) & 23 (25.6\%) & 30 (33.3\%) \\
 &  & CAPITAL & 273 & 117 (42.9\%) & 61 (22.3\%) & 90 (33.0\%) \\
\addlinespace[2pt]
 & CORRIENTES & CAPITAL & 90 & 32 (35.6\%) & 22 (24.4\%) & 23 (25.6\%) \\
\addlinespace[2pt]
 & ENTRE RIOS & PARANA & 90 & 29 (32.2\%) & 19 (21.1\%) & 15 (16.7\%) \\
\addlinespace[2pt]
 & LA PAMPA & TOAY & 90 & 33 (36.7\%) & 17 (18.9\%) & 22 (24.4\%) \\
 &  & UTRACAN & 90 & 33 (36.7\%) & 24 (26.7\%) & 17 (18.9\%) \\
 &  & CAPITAL & 90 & 24 (26.7\%) & 15 (16.7\%) & 15 (16.7\%) \\
\addlinespace[2pt]
 & MENDOZA & CAPITAL & 90 & 29 (32.2\%) & 14 (15.6\%) & 26 (28.9\%) \\
 &  & TUPUNGATO & 182 & 56 (30.8\%) & 32 (17.6\%) & 40 (22.0\%) \\
 &  & GUAYMALLEN & 92 & 28 (30.4\%) & 18 (19.6\%) & 23 (25.0\%) \\
 &  & GODOY CRUZ & 90 & 20 (22.2\%) & 10 (11.1\%) & 17 (18.9\%) \\
\addlinespace[2pt]
 & NEUQUEN & CONFLUENCIA & 91 & 32 (35.2\%) & 17 (18.7\%) & 24 (26.4\%) \\
\addlinespace[2pt]
 & SAN JUAN & RIVADAVIA & 92 & 20 (21.7\%) & 8 (8.7\%) & 15 (16.3\%) \\
\addlinespace[2pt]
 & SANTA FE & GENERAL LOPEZ & 90 & 36 (40.0\%) & 18 (20.0\%) & 26 (28.9\%) \\
 &  & CASTELLANOS & 181 & 72 (39.8\%) & 38 (21.0\%) & 51 (28.2\%) \\
 &  & ROSARIO & 364 & 119 (32.7\%) & 62 (17.0\%) & 104 (28.6\%) \\
\addlinespace[2pt]
 & SANTIAGO DEL ESTERO & GENERAL TABOADA & 180 & 56 (31.1\%) & 18 (10.0\%) & 46 (25.6\%) \\
 &  & ROBLES & 75 & 23 (30.7\%) & 8 (10.7\%) & 21 (28.0\%) \\
 &  & BANDA & 18 & 3 (16.7\%) & 1 (5.6\%) & 2 (11.1\%) \\
\addlinespace[2pt]
 & TUCUMAN & CAPITAL & 92 & 36 (39.1\%) & 23 (25.0\%) & 25 (27.2\%) \\
 &  & TAFI VIAJO & 181 & 58 (32.0\%) & 40 (22.1\%) & 34 (18.8\%) \\
 &  & CRUZ ALTA & 91 & 25 (27.5\%) & 17 (18.7\%) & 16 (17.6\%) \\
 &  & YERBA BUENA & 94 & 18 (19.1\%) & 10 (10.6\%) & 11 (11.7\%) \\
\end{longtable}
\begin{minipage}{0.98\linewidth}
\footnotesize\setlength{\parindent}{0pt}
\textit{Notes:} The table lists municipalities with at least one observation containing a missing value in the full individual-characteristics specification. Each entry reports the number of incomplete observations and, in parentheses, their unweighted percentage among LAPOP observations in that municipality. Full includes age, male, unemployment status, partnership status, high-school completion, left--right ideology, political interest, and blank/null vote, as well as the dependent variable, survey wave, municipality identifier, and survey weight. The other columns apply the corresponding reduced specifications used in Table A10. Municipalities are grouped by province and broad geographic region.
\end{minipage}
\endgroup
\end{landscape}
